\documentclass[11pt]{article}


\usepackage{fullpage}
\usepackage{graphicx}
\usepackage{amsmath}
\usepackage{amssymb}
\usepackage{amsthm}
\usepackage{fancyvrb}

\newcommand{\myname}{Mehshan Mustafa}

\newenvironment{theorem}[2][Theorem]{\begin{trivlist}
\item[\hskip \labelsep {\bfseries #1}\hskip \labelsep {\bfseries #2.}]}{\end{trivlist}}
\newenvironment{lemma}[2][Lemma]{\begin{trivlist}
\item[\hskip \labelsep {\bfseries #1}\hskip \labelsep {\bfseries #2.}]}{\end{trivlist}}
\newenvironment{exercise}[2][Exercise]{\begin{trivlist}
\item[\hskip \labelsep {\bfseries #1}\hskip \labelsep {\bfseries #2.}]}{\end{trivlist}}
\newenvironment{problem}[2][Problem]{\begin{trivlist}
\item[\hskip \labelsep {\bfseries #1}\hskip \labelsep {\bfseries #2.}]}{\end{trivlist}}
\newenvironment{question}[2][Question]{\begin{trivlist}
\item[\hskip \labelsep {\bfseries #1}\hskip \labelsep {\bfseries #2.}]}{\end{trivlist}}
\newenvironment{corollary}[2][Corollary]{\begin{trivlist}
\item[\hskip \labelsep {\bfseries #1}\hskip \labelsep {\bfseries #2.}]}{\end{trivlist}}
\newenvironment{solution}{\begin{proof}[Solution]}{\end{proof}}
\newenvironment{idea}[2][Proof Idea.]{\textit{#1} #2}



\parindent0in
\pagestyle{plain}
\thispagestyle{plain}

\usepackage{csquotes}
\usepackage[shortlabels]{enumitem}

\newcommand{\dated}{\today}

\begin{document}

\textbf{Introduction to the Theory of
Computation}\hfill\textbf{\myname}\\[0.01in]
\textbf{Chapter 1: Regular Languages}\hfill\textbf{\dated}\\
\smallskip\hrule\bigskip

\begin{problem}{1.61}
Consider the languages $C_k$ defined in Problem 1.60. Prove that for each $k$, no DFA
can recognize $C_k$ with fewer than $2^k$ states.
\end{problem}

\begin{proof}
To prove that for each $k$, no DFA can recognize $C_k$ with fewer than $2^k$ states, we give a general solution for constructing the index\footnote{\textbf{Myhill–Nerode theorem}. Let $L$ be a language and let $X$ be a set of strings. Say that $X$ is pairwise distinguishable by $L$ if every two distinct strings in $X$ are distinguishable by $L$. Define the \textbf{\textit{index of L}} to be the maximum number of elements in any set that is pairwise distinguishable by $L$. The index of L may be finite or infinite.} $X_k$ of language $C_k$, where $|X_k| = 2^k$. Construct the index of $C_k$ as follows.
\[
X_1 = \{a,b\}.
\]
\[
X_k = X_{k-1} \bigcup a\Sigma^{k-1}.
\]
For example, $X_2 = \{a,b,aa,ab\}$, and $X_3 = \{a,b,aa,ab,aaa,aab,aba,abb\}$.

% TODO: Prove construction of index works and its size is 2^k.

\end{proof}

\end{document}